
\documentclass{article}
\usepackage{colortbl}
\usepackage{makecell}
\usepackage{multirow}
\usepackage{supertabular}

\begin{document}

\newcounter{utterance}

\centering \large Interaction Transcript for game `dond', experiment `coop\_it', episode 2 with gpt{-}5{-}mini{-}2025{-}08{-}07{-}no{-}reasoning{-}t1.0.
\vspace{24pt}

{ \footnotesize  \setcounter{utterance}{1}
\setlength{\tabcolsep}{0pt}
\begin{supertabular}{c@{$\;$}|p{.15\linewidth}@{}p{.15\linewidth}p{.15\linewidth}p{.15\linewidth}p{.15\linewidth}p{.15\linewidth}}
    \# & $\;$A & \multicolumn{4}{c}{Game Master} & $\;\:$B\\
    \hline

    \theutterance \stepcounter{utterance}  
    & & \multicolumn{4}{p{0.6\linewidth}}{
        \cellcolor[rgb]{0.9,0.9,0.9}{
            \makecell[{{p{\linewidth}}}]{
                \texttt{\tiny{[P1$\langle$GM]}}
                \texttt{State giocando a un gioco di negoziazione cooperativa in cui dovete accordarvi su come dividere una serie di oggetti tra voi e un altro giocatore.} \\
\\ 
\texttt{Regole:} \\
\texttt{(a) A Lei e all'altro giocatore viene dato un insieme di oggetti. Ognuno di voi riceve anche una funzione di valore segreta, che rappresenta il valore di ciascun tipo di oggetto.} \\
\texttt{(b) Si scambiano messaggi con l'altro giocatore per concordare chi si aggiudica gli oggetti. Potete inviare un massimo di 5 messaggi ciascuno, oppure terminare in anticipo facendo una proposta segreta in qualsiasi momento.} \\
\texttt{(c) A ciascuno di voi viene chiesto di inviare una proposta segreta indicando gli oggetti che desiderate, formattata tra parentesi quadre come segue: "{[}Proposta: <numero> <nome oggetto>, <numero> <nome oggetto>, <...>{]}".} \\
\texttt{(d) Se le vostre proposte sono complementari, cioè ci sono abbastanza oggetti per soddisfare entrambe le proposte, a ciascun giocatore viene assegnato un punteggio basato sulla somma dei valori degli oggetti ricevuti. In caso contrario, entrambi ricevono zero punti.} \\
\texttt{(e) L'obiettivo è massimizzare la somma del proprio punteggio e di quello dell'altro giocatore.} \\
\\ 
\texttt{Cominciamo.} \\
\\ 
\texttt{L'insieme degli oggetti disponibili è:} \\
\\ 
\texttt{1 telefono, 2 specchi, 2 frutti} \\
\\ 
\texttt{La funzione valore segreta è:} \\
\\ 
\texttt{telefono: 4, specchio: 2, frutta: 1} \\
\\ 
\texttt{IMPORTANTE: i vostri messaggi, a meno che non si tratti di una proposta segreta, vengono trasmessi direttamente all'altro giocatore, quindi non includete alcuna risposta alle regole o testo di annuncio del vostro messaggio. Per fare una proposta segreta, utilizzate il formato indicato. Non utilizzare delle parentesi quadre quando si comunica all'altro giocatore, altrimenti verrà interpretata come una proposta segreta.} \\
            }
        }
    }
    & \\ \\

    \theutterance \stepcounter{utterance}  
    & \multicolumn{4}{p{0.6\linewidth}}{
        \cellcolor[rgb]{0.9,0.9,0.9}{
            \makecell[{{p{\linewidth}}}]{
                \texttt{\tiny{[P1$\rangle$GM]}}
                \texttt{Ciao! Per me il telefono vale di gran lunga più del resto (4 punti), poi gli specchi (2 ciascuno) e infine la frutta (1 ciascuna). Per massimizzare la somma totale, ha senso che il telefono vada a chi lo valuta di più. A te quanto valgono questi oggetti? In particolare: quanto vale per te il telefono rispetto a specchi e frutta?} \\
            }
        }
    }
    & & \\ \\

    \theutterance \stepcounter{utterance}  
    & & & \multicolumn{2}{p{0.3\linewidth}}{
        \cellcolor[rgb]{0.9,0.9,0.9}{
            \makecell[{{p{\linewidth}}}]{
                \texttt{\tiny{[GM$|$GM]}}
                \texttt{continue} \\
            }
        }
    }
    & & \\ \\

    \theutterance \stepcounter{utterance}  
    & & & \multicolumn{2}{p{0.3\linewidth}}{
        \cellcolor[rgb]{0.9,0.9,0.9}{
            \makecell[{{p{\linewidth}}}]{
                \texttt{\tiny{[GM$|$GM]}}
                \texttt{Ciao! Per me il telefono vale di gran lunga più del resto (4 punti), poi gli specchi (2 ciascuno) e infine la frutta (1 ciascuna). Per massimizzare la somma totale, ha senso che il telefono vada a chi lo valuta di più. A te quanto valgono questi oggetti? In particolare: quanto vale per te il telefono rispetto a specchi e frutta?} \\
            }
        }
    }
    & & \\ \\

    \theutterance \stepcounter{utterance}  
    & & \multicolumn{4}{p{0.6\linewidth}}{
        \cellcolor[rgb]{0.9,0.9,0.9}{
            \makecell[{{p{\linewidth}}}]{
                \texttt{\tiny{[GM$\rangle$P2]}}
                \texttt{State giocando a un gioco di negoziazione cooperativa in cui dovete accordarvi su come dividere una serie di oggetti tra voi e un altro giocatore.} \\
\\ 
\texttt{Regole:} \\
\texttt{(a) A Lei e all'altro giocatore viene dato un insieme di oggetti. Ognuno di voi riceve anche una funzione di valore segreta, che rappresenta il valore di ciascun tipo di oggetto.} \\
\texttt{(b) Si scambiano messaggi con l'altro giocatore per concordare chi si aggiudica gli oggetti. Potete inviare un massimo di 5 messaggi ciascuno, oppure terminare in anticipo facendo una proposta segreta in qualsiasi momento.} \\
\texttt{(c) A ciascuno di voi viene chiesto di inviare una proposta segreta indicando gli oggetti che desiderate, formattata tra parentesi quadre come segue: "{[}Proposta: <numero> <nome oggetto>, <numero> <nome oggetto>, <...>{]}".} \\
\texttt{(d) Se le vostre proposte sono complementari, cioè ci sono abbastanza oggetti per soddisfare entrambe le proposte, a ciascun giocatore viene assegnato un punteggio basato sulla somma dei valori degli oggetti ricevuti. In caso contrario, entrambi ricevono zero punti.} \\
\texttt{(e) L'obiettivo è massimizzare la somma del proprio punteggio e di quello dell'altro giocatore.} \\
\\ 
\texttt{Cominciamo.} \\
\\ 
\texttt{L'insieme degli oggetti disponibili è:} \\
\\ 
\texttt{1 telefono, 2 specchi, 2 frutti} \\
\\ 
\texttt{La funzione valore segreta è:} \\
\\ 
\texttt{telefono: 6, specchio: 0, frutta: 2} \\
\\ 
\texttt{IMPORTANTE: i vostri messaggi, a meno che non si tratti di una proposta segreta, vengono trasmessi direttamente all'altro giocatore, quindi non includete alcuna risposta alle regole o testo di annuncio del vostro messaggio. Per fare una proposta segreta, utilizzate il formato indicato. Non utilizzare delle parentesi quadre quando si comunica all'altro giocatore, altrimenti verrà interpretata come una proposta segreta.} \\
\\ 
\\ 
\texttt{Ciao! Per me il telefono vale di gran lunga più del resto (4 punti), poi gli specchi (2 ciascuno) e infine la frutta (1 ciascuna). Per massimizzare la somma totale, ha senso che il telefono vada a chi lo valuta di più. A te quanto valgono questi oggetti? In particolare: quanto vale per te il telefono rispetto a specchi e frutta?} \\
            }
        }
    }
    & \\ \\

    \theutterance \stepcounter{utterance}  
    & & & \multicolumn{4}{p{0.6\linewidth}}{
        \cellcolor[rgb]{0.9,0.9,0.9}{
            \makecell[{{p{\linewidth}}}]{
                \texttt{\tiny{[GM$\langle$P2]}}
                \texttt{Grazie della trasparenza! Per me i valori sono: telefono 6, specchio 0, frutta 2 ciascuno. Quindi anch’io valuto il telefono più di tutto, poi la frutta, mentre gli specchi per me valgono zero.} \\
\\ 
\texttt{Dato che tu valuti il telefono 4 (più di ogni altro singolo oggetto), e io 6, la somma totale è massima se il telefono va a me (6 invece di 4). Per massimizzare anche la tua utilità, potremmo compensare con gli specchi (che per me valgono 0 ma per te 2 ciascuno) e parte della frutta.} \\
\\ 
\texttt{Proposta aperta: io prendo il telefono e 1 frutto; tu prendi 2 specchi e 1 frutto.} \\
\texttt{{-} Mio valore: 6 + 2 = 8} \\
\texttt{{-} Tuo valore: 2+2 + 1 = 5} \\
\texttt{Somma totale: 13 (massima possibile: 6 + 2 + 2 + 1 + 1 = 12? No, ricontrollo: la somma massima dei valori è 6 (tel) + 0+0 (specchi per me) + 2+2 (frutti per me) = 10 se vanno tutti a me; ma considerando i tuoi valori sugli specchi e la frutta, distribuendo bene possiamo superare 10. Con la divisione proposta otteniamo 8+5=13, che è meglio per la somma).} \\
\\ 
\texttt{Se ti va bene, conferma; altrimenti proponi una variante (ad es. telefono a me, a te 2 specchi e 2 frutti, e a me 0 frutti), così vediamo l’equilibrio che preferisci.} \\
            }
        }
    }
     \\ \\

    \theutterance \stepcounter{utterance}  
    & & & \multicolumn{2}{p{0.3\linewidth}}{
        \cellcolor[rgb]{0.9,0.9,0.9}{
            \makecell[{{p{\linewidth}}}]{
                \texttt{\tiny{[GM$|$GM]}}
                \texttt{continue} \\
            }
        }
    }
    & & \\ \\

    \theutterance \stepcounter{utterance}  
    & & & \multicolumn{2}{p{0.3\linewidth}}{
        \cellcolor[rgb]{0.9,0.9,0.9}{
            \makecell[{{p{\linewidth}}}]{
                \texttt{\tiny{[GM$|$GM]}}
                \texttt{Grazie della trasparenza! Per me i valori sono: telefono 6, specchio 0, frutta 2 ciascuno. Quindi anch’io valuto il telefono più di tutto, poi la frutta, mentre gli specchi per me valgono zero.} \\
\\ 
\texttt{Dato che tu valuti il telefono 4 (più di ogni altro singolo oggetto), e io 6, la somma totale è massima se il telefono va a me (6 invece di 4). Per massimizzare anche la tua utilità, potremmo compensare con gli specchi (che per me valgono 0 ma per te 2 ciascuno) e parte della frutta.} \\
\\ 
\texttt{Proposta aperta: io prendo il telefono e 1 frutto; tu prendi 2 specchi e 1 frutto.} \\
\texttt{{-} Mio valore: 6 + 2 = 8} \\
\texttt{{-} Tuo valore: 2+2 + 1 = 5} \\
\texttt{Somma totale: 13 (massima possibile: 6 + 2 + 2 + 1 + 1 = 12? No, ricontrollo: la somma massima dei valori è 6 (tel) + 0+0 (specchi per me) + 2+2 (frutti per me) = 10 se vanno tutti a me; ma considerando i tuoi valori sugli specchi e la frutta, distribuendo bene possiamo superare 10. Con la divisione proposta otteniamo 8+5=13, che è meglio per la somma).} \\
\\ 
\texttt{Se ti va bene, conferma; altrimenti proponi una variante (ad es. telefono a me, a te 2 specchi e 2 frutti, e a me 0 frutti), così vediamo l’equilibrio che preferisci.} \\
            }
        }
    }
    & & \\ \\

    \theutterance \stepcounter{utterance}  
    & & \multicolumn{4}{p{0.6\linewidth}}{
        \cellcolor[rgb]{0.9,0.9,0.9}{
            \makecell[{{p{\linewidth}}}]{
                \texttt{\tiny{[P1$\langle$GM]}}
                \texttt{Grazie della trasparenza! Per me i valori sono: telefono 6, specchio 0, frutta 2 ciascuno. Quindi anch’io valuto il telefono più di tutto, poi la frutta, mentre gli specchi per me valgono zero.} \\
\\ 
\texttt{Dato che tu valuti il telefono 4 (più di ogni altro singolo oggetto), e io 6, la somma totale è massima se il telefono va a me (6 invece di 4). Per massimizzare anche la tua utilità, potremmo compensare con gli specchi (che per me valgono 0 ma per te 2 ciascuno) e parte della frutta.} \\
\\ 
\texttt{Proposta aperta: io prendo il telefono e 1 frutto; tu prendi 2 specchi e 1 frutto.} \\
\texttt{{-} Mio valore: 6 + 2 = 8} \\
\texttt{{-} Tuo valore: 2+2 + 1 = 5} \\
\texttt{Somma totale: 13 (massima possibile: 6 + 2 + 2 + 1 + 1 = 12? No, ricontrollo: la somma massima dei valori è 6 (tel) + 0+0 (specchi per me) + 2+2 (frutti per me) = 10 se vanno tutti a me; ma considerando i tuoi valori sugli specchi e la frutta, distribuendo bene possiamo superare 10. Con la divisione proposta otteniamo 8+5=13, che è meglio per la somma).} \\
\\ 
\texttt{Se ti va bene, conferma; altrimenti proponi una variante (ad es. telefono a me, a te 2 specchi e 2 frutti, e a me 0 frutti), così vediamo l’equilibrio che preferisci.} \\
            }
        }
    }
    & \\ \\

    \theutterance \stepcounter{utterance}  
    & \multicolumn{4}{p{0.6\linewidth}}{
        \cellcolor[rgb]{0.9,0.9,0.9}{
            \makecell[{{p{\linewidth}}}]{
                \texttt{\tiny{[P1$\rangle$GM]}}
                \texttt{La tua proposta è ottima per la somma totale e per me è già molto buona. Confermo che mi va bene: telefono a te, 2 specchi a me, 1 frutto a te e 1 a me.} \\
\\ 
\texttt{Per chiudere ufficialmente, invio la proposta segreta ora.} \\
\\ 
\texttt{{[}Proposta: 2 specchi, 1 frutta{]}} \\
            }
        }
    }
    & & \\ \\

    \theutterance \stepcounter{utterance}  
    & & & \multicolumn{2}{p{0.3\linewidth}}{
        \cellcolor[rgb]{0.9,0.9,0.9}{
            \makecell[{{p{\linewidth}}}]{
                \texttt{\tiny{[GM$|$GM]}}
                \texttt{proposal} \\
            }
        }
    }
    & & \\ \\

    \theutterance \stepcounter{utterance}  
    & & & \multicolumn{2}{p{0.3\linewidth}}{
        \cellcolor[rgb]{0.9,0.9,0.9}{
            \makecell[{{p{\linewidth}}}]{
                \texttt{\tiny{[GM$|$GM]}}
                [0, 2, 1]
            }
        }
    }
    & & \\ \\

    \theutterance \stepcounter{utterance}  
    & & \multicolumn{4}{p{0.6\linewidth}}{
        \cellcolor[rgb]{0.9,0.9,0.9}{
            \makecell[{{p{\linewidth}}}]{
                \texttt{\tiny{[GM$\rangle$P2]}}
                \texttt{L'altro giocatore ha presentato una proposta segreta. Ora tocca a Lei presentare una sua proposta.} \\
            }
        }
    }
    & \\ \\

    \theutterance \stepcounter{utterance}  
    & & & \multicolumn{4}{p{0.6\linewidth}}{
        \cellcolor[rgb]{0.9,0.9,0.9}{
            \makecell[{{p{\linewidth}}}]{
                \texttt{\tiny{[GM$\langle$P2]}}
                \texttt{{[}Proposta: 1 telefono, 2 specchi, 1 frutto{]}} \\
            }
        }
    }
     \\ \\

    \theutterance \stepcounter{utterance}  
    & & & \multicolumn{2}{p{0.3\linewidth}}{
        \cellcolor[rgb]{0.9,0.9,0.9}{
            \makecell[{{p{\linewidth}}}]{
                \texttt{\tiny{[GM$|$GM]}}
                \texttt{proposal} \\
            }
        }
    }
    & & \\ \\

    \theutterance \stepcounter{utterance}  
    & & & \multicolumn{2}{p{0.3\linewidth}}{
        \cellcolor[rgb]{0.9,0.9,0.9}{
            \makecell[{{p{\linewidth}}}]{
                \texttt{\tiny{[GM$|$GM]}}
                [1, 2, 1]
            }
        }
    }
    & & \\ \\

    \theutterance \stepcounter{utterance}  
    & & & \multicolumn{2}{p{0.3\linewidth}}{
        \cellcolor[rgb]{0.9,0.9,0.9}{
            \makecell[{{p{\linewidth}}}]{
                \texttt{\tiny{[GM$|$GM]}}
                [[0, 2, 1], [1, 2, 1]]
            }
        }
    }
    & & \\ \\

\end{supertabular}
}

\end{document}
